% ============================================================
% SWA-MPPI v3: Cross-Cultural Moral Alignment of LLMs
% NeurIPS 2026 Submission -- REVISED VERSION
% Compile: pdflatex -> bibtex -> pdflatex -> pdflatex
% ============================================================
\documentclass{article}
\usepackage[final,main]{neurips_2026}

\usepackage[utf8]{inputenc}
\usepackage[T1]{fontenc}
\usepackage{hyperref}
\usepackage{url}
\usepackage{booktabs}
\usepackage{amsfonts}
\usepackage{amsmath}
\usepackage{amssymb}
\usepackage{nicefrac}
\usepackage{microtype}
\usepackage{xcolor}
\usepackage{graphicx}
\usepackage{subcaption}
\usepackage{algorithm}
\usepackage{algpseudocode}
\usepackage{multirow}
\usepackage{makecell}
\usepackage{bbm}

% Colour macros
\definecolor{improvegreen}{RGB}{0,128,0}
\definecolor{worsered}{RGB}{180,0,0}
\newcommand{\gain}[1]{\textcolor{improvegreen}{#1}}
\newcommand{\loss}[1]{\textcolor{worsered}{#1}}

% ============================================================
\title{%
  \textbf{SWA-MPPI}: Socially-Weighted Alignment with\\
  Model Predictive Path Integral Control\\
  for Cross-Cultural Moral Alignment of Large Language Models%
}

\author{%
  Anonymous Authors\\
  \texttt{anonymous@institution.edu}
}

\begin{document}

\maketitle

% ============================================================
\begin{abstract}
% ============================================================
Large language models inherit the moral intuitions of their predominantly Western
training data, yet serve culturally diverse populations worldwide. No existing
method can adapt an LLM's moral reasoning to a target culture \emph{at inference
time} without labelled preference data or weight modification. We introduce
\textbf{SWA-MPPI}, the first training-free framework for cross-cultural moral
alignment that operates entirely in logit space. SWA-MPPI constructs an ensemble
of culturally grounded persona agents from World Values Survey data, detects
inter-persona disagreement as a signal of cultural contestation, and resolves it
via a Prospect-Theory-weighted Model Predictive Path Integral (MPPI) update---a
principled stochastic optimisation that balances private cultural utility with
collective welfare. On the Multilingual Trolley Problem benchmark across 15
countries, 12 native languages, and 6 moral dimensions, SWA-MPPI reduces
Jensen-Shannon Distance to human moral preferences by \textbf{37.6\%}
(Qwen2.5-72B) and \textbf{34.7\%} (Llama-3.1-70B), outperforming all prompting
and activation-steering baselines by $\geq$10\% while operating on frozen weights
with a single additional batched forward pass. Notably, the largest gains occur
for Anglosphere countries ($+$72--75\%), revealing that RLHF-induced biases
actively worsen alignment even for the best-represented populations. The framework
is self-limiting: when the model is already well-calibrated, inter-persona
consensus is high and no correction is applied. Code available at
\texttt{[repository]}.
\end{abstract}

% ============================================================
\section{Introduction}
\label{sec:intro}
% ============================================================

When a self-driving car must choose between sparing a group of elderly men or a
group of young women, there is no culturally universal answer. The Moral Machine
experiment~\citep{awad2018moral}---40~million judgments from 233~countries---revealed
that such preferences vary profoundly across cultures: East Asian respondents weight
youth more heavily; Southern-hemisphere populations show distinct utilitarian calculi;
Islamic societies prioritise different social roles. Yet the large language models
increasingly enlisted to assist with such decisions encode a single, culturally
monolithic moral stance inherited from their predominantly Western training
data~\citep{henrich2010weirdest}.

The multilingual successor benchmark MultiTP~\citep{jin2025multitp} operationalises
this evidence into country-specific Average Marginal Component Effects (AMCE)
across six moral dimensions and 107 languages, providing a rigorous test bed for
cross-cultural alignment. The verdict is stark: \citet{jin2025multitp} found that
none of 19~contemporary LLMs achieved broadly satisfactory alignment---most exhibit
WEIRD-biased reasoning that departs markedly from non-WEIRD human preferences.

How should this gap be closed? Existing approaches each sacrifice critical properties.
\emph{Fine-tuning}~\citep{bai2022constitutional, ouyang2022instructgpt} demands
per-culture labelled data and modifies weights irreversibly.
\emph{Prompting-based steering} relies on ad-hoc instructions with no empirical
grounding~\citep{jin2025multitp}.
\emph{Activation steering}~\citep{arditi2025refusal, zou2023representation} requires
white-box access and per-model engineering.
\emph{Reward-guided decoding}~\citep{khanov2024args, mudgal2024controlled} needs trained
reward models unavailable for cross-cultural settings.
\emph{Culture-aware adapters}~\citep{li2024culturemanager, yao2024caredio} scale poorly
across 100+ countries. None of these methods can adapt an LLM's moral reasoning to a
target culture \emph{at inference time} without labelled data or weight modification.

We propose \textbf{SWA-MPPI} (Socially-Weighted Alignment with Model Predictive Path
Integral control), built on a simple insight: \emph{disagreement is signal}.
Moral preferences within any society are not monolithic---age cohorts and social
strata exhibit systematically different priorities~\citep{awad2018moral,
kleiman2017learning}. The World Values Survey~\citep{haerpfer2022wvs} provides
empirically calibrated data on these within-country differences across 80+ nations.
When culturally grounded persona agents derived from WVS data \emph{disagree} on a
moral decision, that disagreement reveals the decision is culturally contested---and
can be exploited through stochastic optimal control~\citep{williams2017mppi} to
produce a more culturally representative output.

Concretely, SWA-MPPI constructs $N=4$ WVS-grounded personas per country (three
age cohorts plus a utilitarian agent), evaluates all $N+1$ agents in a single batched
forward pass with native-language prompting across 12 languages, and extracts
decision logits. If inter-agent variance exceeds a calibrated threshold~$\tau$, it
samples $K=128$ perturbations and scores each via a Prospect Theory value
function~\citep{kahneman1979prospect} that captures the asymmetric gain-loss
sensitivity characteristic of human moral judgment. The importance-weighted MPPI
update yields a minimal-KL correction; a positional debiasing step cancels recency
effects.

\paragraph{Contributions.}
\begin{itemize}
  \item We introduce SWA-MPPI, the first application of MPPI in logit space for
        cross-cultural value alignment: a training-free framework combining
        WVS-grounded personas, adaptive conflict detection, and Prospect-Theory-weighted
        optimisation (\S\ref{sec:method}).
  \item We demonstrate mean JSD reductions of \textbf{37.6\%} (Qwen2.5-72B) and
        \textbf{34.7\%} (Llama-3.1-70B) across 15 countries, with gains reaching
        $+$75\% for Anglosphere pairs, outperforming all baselines by $\geq$10\%
        (\S\ref{sec:results}).
  \item We conduct systematic ablations revealing a clear component hierarchy:
        the utilitarian persona is most critical ($-$39.1\% relative Pearson without it),
        followed by MPPI ($+$15.9\% JSD), debiasing ($+$12.0\%), and Prospect Theory
        ($+$4.5\%) (\S\ref{sec:ablation}).
\end{itemize}

% ============================================================
\section{Related Work}
\label{sec:related}
% ============================================================

\paragraph{Cross-cultural moral benchmarks.}
The Moral Machine experiment~\citep{awad2018moral} provides 40M+ cross-cultural moral
judgments. MultiTP~\citep{jin2025multitp} extends this to 107 languages with
country-specific AMCEs---we adopt it as our benchmark and propose an alignment
\emph{intervention} rather than diagnostics alone. Other benchmarks
(ETHICS~\citep{hendrycks2021aligning}, MoralExceptQA~\citep{jin2022exceptions},
MoCa~\citep{nie2023moca}) remain English-centric.
PRISM~\citep{kirk2024prism} and GlobalOpinionQA~\citep{durmus2023global} address
diversity but cover few languages. Recent frameworks
(EvalMORAAL~\citep{aakanksha2024evalmor}, CARB~\citep{myung2024carb}) emphasise
broader-topic evaluations; SWA-MPPI targets binary moral decisions specifically.

\paragraph{Cross-cultural LLM evaluation.}
LLM outputs correlate most strongly with liberal, educated American
respondents~\citep{santurkar2023opinions}. RLHF introduces unintended cultural
impacts~\citep{ryan2024unintended}, and assigning cultural personas can surface
toxic outputs~\citep{deshpande2023toxicity}. This recurring WEIRD bias motivates
SWA-MPPI's calibration objective.

\paragraph{Inference-time alignment.}
Activation steering~\citep{arditi2025refusal, zou2023representation} requires
white-box access. ARGS~\citep{khanov2024args} and controlled
decoding~\citep{mudgal2024controlled} need trained reward models---unavailable for
cross-cultural settings (Appendix~\ref{app:baselines}).
CultureManager~\citep{li2024culturemanager} and CAReDiO~\citep{yao2024caredio}
train culture-specific adapters but require per-culture labelled data and weight
modification, scaling poorly across 100+ countries. SWA-MPPI is the first
\emph{training-free} method addressing multi-cultural calibration.
\citet{sorensen2024roadmap} articulate the goal of pluralistic alignment, which
SWA-MPPI operationalises concretely.

\paragraph{MPPI and Prospect Theory.}
MPPI~\citep{williams2017mppi} is a sampling-based stochastic optimal control algorithm
originally developed for robotics; \citet{huang2024mppi} adapted it for token-level
LLM steering. We are the first to apply MPPI in \emph{logit space} for moral
alignment, combining it with Prospect Theory~\citep{kahneman1979prospect} to capture
asymmetric gain-loss sensitivity in human moral judgment.

% ============================================================
\section{Method: SWA-MPPI}
\label{sec:method}
% ============================================================

\subsection{Problem Formulation}
\label{sec:method:problem}

Let $\mathcal{X}$ be the set of moral dilemma scenarios, each presenting a forced
binary choice via decision tokens \texttt{LEFT}~($\ell$) and \texttt{RIGHT}~($r$).
For target country $c$, let $\mathbf{h}^{(c)} \in \mathbb{R}^6$ be the human AMCE
vector from 40M+ judgments~\citep{awad2018moral}. A frozen LLM $f_\theta$ produces
logits $\mathbf{z}(\mathbf{x}) = [z_\ell,\, z_r]$; the spare probability is
$p_{\text{spare}} = \sigma(\Delta / T_{\text{dec}})$ where $\Delta = z_r - z_\ell$.
Each dimension's AMCE is estimated as:
\begin{equation}
  \hat{m}_d^{(c)} = \frac{1}{|\mathcal{S}_d^{(c)}|}\sum_{\mathbf{x} \in \mathcal{S}_d^{(c)}}
  p_{\text{spare}}(\mathbf{x}),
  \label{eq:mpr}
\end{equation}
following the MultiTP evaluation convention. For the five binary dimensions
(Species, Gender, Age, Fitness, Social Value), this reduces to the empirical
mean of $p_{\text{spare}}$. For the Utilitarianism dimension, where the
predictor is the \emph{count difference} $|n_{\text{pref}} - n_{\text{non-pref}}|$
between groups, we fit $p_{\text{spare}} \sim a + b \cdot n_{\text{diff}}$ via
OLS and evaluate at the mean observed $n_{\text{diff}}$
(Appendix~\ref{app:amce}).

\textbf{Goal:} Without modifying $\theta$, find a logit correction
$\delta^*(\mathbf{x}, c)$ such that the corrected AMCE vector
$\hat{\mathbf{m}}^{(c)}(\theta, \delta^*)$ is closer to $\mathbf{h}^{(c)}$
under Jensen-Shannon Divergence.

\subsection{Cultural Persona Construction}
\label{sec:method:personas}

For each country $c$, SWA-MPPI constructs $N=4$ persona agents from WVS Wave~7
individual-level microdata~\citep{haerpfer2022wvs}. We extract 8 value dimensions
from specific WVS items: \emph{gender egalitarianism} (Q58P, Q59P, Q60P),
\emph{religious importance} (Q6P), \emph{interpersonal trust} (Q43P),
\emph{moral permissiveness} (Q50, Q52P, Q54P), \emph{work centrality} (Q5P),
\emph{family importance} (Q1P), \emph{personal autonomy} (Q39P), and
\emph{meritocratic orientation} (Q40P). For each country, respondents are
segmented into three age cohorts using birth year and survey year:
\emph{young} (18--35), \emph{middle-aged} (36--55), and \emph{older} ($>$55).
Dimension means are computed per cohort and converted into natural-language
descriptors (e.g., a mean religiosity ratio $>$0.85 yields ``deeply religious'').
A fourth, country-invariant \emph{utilitarian agent} maximises lives saved.
Each persona is a natural-language system prompt:

\begin{quote}
\small\itshape
You are a [age-cohort] adult from [country]. The cultural values typical of your
social group are: [WVS statistics]. When making moral judgments, you tend to [value
description reflecting WVS profile].
\end{quote}

\noindent WVS grounding ensures empirically measured within-country variation rather
than stereotypes~\citep{deshpande2023toxicity}. Concrete persona examples and the
WVS-to-trolley dimension linkage are detailed in Appendix~\ref{app:personas}.

\subsection{Batched Evaluation and Logit Extraction}
\label{sec:method:eval}

All $N+1=5$ input sequences (base + 4 personas) are evaluated in a \emph{single}
batched forward pass:
\begin{equation}
  \mathbf{z}_{\text{base}},\; \mathbf{z}_1, \ldots, \mathbf{z}_N
  = f_\theta\!\left(\mathbf{x};\; \varnothing,\; \mathbf{s}_1, \ldots, \mathbf{s}_N\right).
  \label{eq:batch}
\end{equation}
Per-category logit temperatures $T_{\text{cat}}$ (Table~\ref{tab:hyperparams}) scale
logits before computing decision gaps:
\begin{align}
  \delta_i &= \frac{z_{i,r} - z_{i,\ell}}{T_{\text{cat}}}, \quad
  \bar{\delta} = \frac{1}{N}\sum_{i=1}^{N} \delta_i \quad \text{(consensus)},
  \label{eq:consensus} \\
  r_i &= \delta_i - \delta_{\text{base}} \quad \text{(per-agent reward)}.
  \label{eq:reward}
\end{align}
The reward $r_i$ measures how strongly persona $i$ shifts the model \emph{away from
the culturally-neutral baseline}---a large $|r_i|$ indicates that the persona has
identified a direction of moral preference the vanilla model does not exhibit.

\subsection{Adaptive Conflict Detection}
\label{sec:method:tau}

When personas agree, the consensus $\bar{\delta}$ suffices. When they disagree
substantially, a controlled correction is warranted. SWA-MPPI uses a per-country
threshold $\tau^{(c)}$ calibrated on $n_{\text{cal}}=50$ unlabelled scenarios to
target a 35\% MPPI trigger rate:
\begin{equation}
  \tau^{(c)} = \operatorname{Percentile}_{65}
  \!\Bigl(\bigl\{\operatorname{Var}(\delta_1^{(s)}, \ldots, \delta_N^{(s)})\bigr\}_{s=1}^{n_{\text{cal}}}\Bigr).
  \label{eq:tau}
\end{equation}
MPPI fires iff $\operatorname{Var}(\delta_1, \ldots, \delta_N) \geq \tau^{(c)}$.
Calibration uses only variance statistics and never observes human AMCEs
(Appendix~\ref{app:calibration}).

\subsection{Prospect-Theory-Weighted MPPI Optimisation}
\label{sec:method:mppi}

When triggered, SWA-MPPI draws $K=128$ perturbations
$\epsilon_k \sim \mathcal{N}(0,\sigma^2)$ ($\sigma=0.3$) and evaluates a collective
welfare score for each $\tilde{\delta}_k = \bar{\delta} + \epsilon_k$.

\paragraph{Per-agent cooperative utility.}
\begin{equation}
  U_i(\epsilon_k) = (1 - \lambda_{\text{coop}})\,v(r_i\,\tilde{\delta}_k)
  \;+\; \lambda_{\text{coop}}\,v\!\Bigl(\bar{r}_{-i}\,\tilde{\delta}_k\Bigr),
  \label{eq:util_agent}
\end{equation}
where $\lambda_{\text{coop}}=0.7$ and $v(\cdot)$ is the Prospect Theory value
function with $\alpha=\beta=0.88$, $\kappa=2.25$~\citep{kahneman1979prospect}:
\begin{equation}
  v(x) = \begin{cases}
    x^{\alpha} & x \geq 0,\\
    -\kappa\,|x|^{\beta} & x < 0.
  \end{cases}
  \label{eq:prospect}
\end{equation}

\paragraph{Collective utility with KL regularisation.}
\begin{equation}
  U_{\text{total}}(\epsilon_k)
  = \frac{1}{N}\sum_{i=1}^{N} U_i(\epsilon_k)
  \;-\; \alpha_{\text{KL}}\,\frac{\epsilon_k^2}{2\sigma^2},
  \label{eq:util_total}
\end{equation}
where $\alpha_{\text{KL}}=0.05$. The KL term derives from the free-energy objective
of MPPI, preventing overcorrection in a theoretically grounded way
(Appendix~\ref{app:mppi_theory}).

\paragraph{MPPI update.}
\begin{equation}
  w_k = \frac{\exp\!\bigl(U_{\text{total}}(\epsilon_k)/\beta\bigr)}
             {\sum_{k'}\exp\!\bigl(U_{\text{total}}(\epsilon_{k'})/\beta\bigr)},
  \qquad
  \delta^* = \sum_{k=1}^{K} w_k\,\epsilon_k,
  \label{eq:mppi}
\end{equation}
where $\beta = T_{\text{dec}} = 0.5$ is the MPPI inverse temperature that governs the
sharpness of importance weights. We set $\beta$ equal to the decision temperature
by design: this ensures that the MPPI optimisation operates on the same scale as
the final sigmoid decision, so that a perturbation $\epsilon_k$ that shifts
$\delta_{\text{opt}}$ by one unit produces a proportionate change in
$p_{\text{spare}}$. Lower $\beta$ would concentrate weights on fewer high-utility
samples (sharper selection); higher $\beta$ would spread weights more uniformly
(softer averaging). At $\beta = 0.5$, the effective sample size
$K_{\text{eff}} = 1/\sum_k w_k^2$ is typically 15--40 (out of $K=128$), indicating
moderate concentration that balances exploration with exploitation.

\subsection{Positional Debiasing and Final Prediction}
\label{sec:method:debias}

To cancel recency bias~\citep{liu2024lost}, each scenario is evaluated twice: once
in the original ordering, once with \emph{both} lane labels (LEFT $\leftrightarrow$
RIGHT) and group labels (Group~A $\leftrightarrow$ Group~B) swapped in the native
language. The debiased decision gap for each agent is
$\delta_i = (\delta_i^{(\text{orig})} - \delta_i^{(\text{swap})})/2$.
The final prediction is $p_{\text{spare}} = \sigma(\delta_{\text{opt}} / T_{\text{dec}})$
where $\delta_{\text{opt}} = \bar{\delta} + \delta^*$ and $T_{\text{dec}} = 0.5$.

Algorithm~\ref{alg:swa} summarises the full procedure; Table~\ref{tab:hyperparams}
lists all hyperparameters (validation details in Appendix~\ref{app:hyperparams}).

\begin{algorithm}[t]
\caption{SWA-MPPI Inference for Scenario $\mathbf{x}$ in Country $c$}
\label{alg:swa}
\begin{algorithmic}[1]
\Require Frozen LLM $f_\theta$; personas $\{\mathbf{s}_i\}_{i=1}^{N}$;
         threshold $\tau^{(c)}$; hyperparameters $K,\sigma,\lambda_{\text{coop}},\alpha_{\text{KL}},T_{\text{dec}}$
\Ensure $p_{\text{spare}} \in [0,1]$
\For{$\text{order} \in \{\text{orig}, \text{swap}\}$}
  \State Batched forward: evaluate $f_\theta$ with base + $N$ personas
  \State Extract logit pairs; compute $\delta_i^{(\text{ord})}$ \hfill\Comment{Eq.~\ref{eq:batch}}
\EndFor
\State Average: $\delta_i \leftarrow (\delta_i^{(\text{orig})} - \delta_i^{(\text{swap})})/2$ \hfill\Comment{Debiasing}
\State Compute consensus $\bar{\delta}$, rewards $r_i$, variance $\nu$ \hfill\Comment{Eqs.~\ref{eq:consensus}--\ref{eq:reward}}
\If{$\nu \geq \tau^{(c)}$} \hfill\Comment{Conflict detected}
  \For{$k = 1$ \textbf{to} $K$}
    \State Sample $\epsilon_k \sim \mathcal{N}(0, \sigma^2)$; compute $U_{\text{total}}(\epsilon_k)$ \hfill\Comment{Eqs.~\ref{eq:util_agent}--\ref{eq:util_total}}
  \EndFor
  \State $\delta^* = \sum_k w_k \epsilon_k$ \hfill\Comment{Eq.~\ref{eq:mppi}}
\Else \quad $\delta^* \leftarrow 0$
\EndIf
\State \Return $p_{\text{spare}} = \sigma((\bar{\delta} + \delta^*) / T_{\text{dec}})$
\end{algorithmic}
\end{algorithm}

\begin{table}[t]
\caption{SWA-MPPI hyperparameters. PT parameters from \citet{kahneman1979prospect};
others validated on a held-out synthetic set (Appendix~\ref{app:hyperparams}).}
\label{tab:hyperparams}
\centering\small
\begin{tabular}{llll}
\toprule
Parameter & Symbol & Value & Source \\
\midrule
Persona agents        & $N$                     & 4     & Design \\
MPPI samples          & $K$                     & 128   & Design \\
Perturbation $\sigma$ & $\sigma$                & 0.3   & Validation \\
Cooperation weight    & $\lambda_{\text{coop}}$ & 0.7   & Validation \\
KL penalty            & $\alpha_{\text{KL}}$    & 0.05  & Validation \\
PT curvature          & $\alpha{=}\beta$        & 0.88  & \citet{kahneman1979prospect} \\
PT loss aversion      & $\kappa$                & 2.25  & \citet{kahneman1979prospect} \\
Decision temp.        & $T_{\text{dec}}$        & 0.5   & Validation \\
Trigger rate          & $\rho_{\text{target}}$  & 0.35  & Design \\
$T_{\text{cat}}$ (Species / Gender / Others) & & 4.0 / 3.5 / 1.5 & Category \\
\bottomrule
\end{tabular}
\end{table}

% -- Architecture Figure --
\begin{figure}[t]
  \centering
  \fbox{\parbox{\dimexpr\linewidth-2\fboxsep-2\fboxrule}{%
    \centering\vspace{3cm}%
    \textbf{Figure placeholder:} SWA-MPPI Pipeline Overview\\[4pt]
    \small Left: WVS Persona Construction \quad$\rightarrow$\quad
    Middle: Batched Forward Pass \quad$\rightarrow$\quad
    Right: MPPI Optimisation Loop%
    \vspace{3cm}%
  }}
  \caption{%
    \textbf{SWA-MPPI pipeline.}
    \textbf{Left:} WVS-grounded persona construction (3 age cohorts + utilitarian).
    \textbf{Middle:} Batched forward pass with positional debiasing.
    \textbf{Right:} MPPI optimisation when inter-persona variance exceeds $\tau^{(c)}$.
  }
  \label{fig:overview}
\end{figure}

% ============================================================
\section{Experimental Setup}
\label{sec:experiments}
% ============================================================

\subsection{Dataset}
\label{sec:exp:data}

We evaluate on MultiTP~\citep{jin2025multitp}: trolley-problem vignettes in native
languages across six moral dimensions (Table~\ref{tab:categories}). For each of 15
countries, we apply: (i)~deduplication (retain only original-phrasing vignettes),
(ii)~quality filtering (remove Utilitarianism scenarios where both groups have equal
size, as these carry no utilitarian signal), (iii)~category capping (max 80 per
dimension), and (iv)~up-sampling via random oversampling with replacement (min 36
per dimension for stable AMCE estimation), yielding 310--500 scenarios per country
(Appendix~\ref{app:dataset}). The adaptive $\tau^{(c)}$ is calibrated on 50
unlabelled scenarios using only inter-agent variance statistics---never human AMCEs.

\begin{table}[t]
\caption{The six moral dimensions in MultiTP.}
\label{tab:categories}
\centering\small
\begin{tabular}{lllc}
\toprule
Dimension      & Preferred       & Contra          & $T_{\text{cat}}$ \\
\midrule
Species        & Human           & Animal          & 4.0 \\
Gender         & Female          & Male            & 3.5 \\
Age            & Young           & Elderly         & 1.5 \\
Fitness        & Fit             & Unfit           & 1.5 \\
Social Value   & High (exec.)    & Low (homeless)  & 1.5 \\
Utilitarianism & More lives      & Fewer lives     & 1.5 \\
\bottomrule
\end{tabular}
\end{table}

\subsection{Models}
\label{sec:exp:models}

Six LLMs: \textbf{Qwen2.5-72B-Instruct}~\citep{qwen2025technical} (Int8),
\textbf{Llama-3.1-70B-Instruct}~\citep{grattafiori2024llama} (FP8),
\textbf{Qwen2.5-7B/32B-Instruct} and \textbf{Llama-3.1-8B-Instruct} (4-bit),
\textbf{Mistral-Large-2407}~\citep{mistral2024} (4-bit). All on a single H100 80GB
via Unsloth~\citep{unsloth2024}, greedy decoding, seed 42.

\textbf{Native-language evaluation.} Each country's scenarios are presented
\emph{entirely in the native language}: scenario framing, character descriptions,
lane labels, and closing questions are all translated into 12 languages (zh, ja, ko,
de, fr, pt, ar, vi, hi, ru, es, en). Only the \emph{decision tokens}
(\texttt{LEFT}/\texttt{RIGHT}) remain English---these are verified as single-token
entries in each model's vocabulary (Appendix~\ref{app:token_elicitation}). This
design ensures that the moral reasoning is grounded in the target language's
conceptual framing while maintaining a consistent, unambiguous decision interface
across all models and languages.

\subsection{Baselines}
\label{sec:exp:baselines}

On Llama-3.1-70B: (1)~\textbf{Vanilla} (generic system prompt),
(2)~\textbf{Country-Tailored Prompt (B1)},
(3)~\textbf{WVS Profile Prompting (B2)},
(4)~\textbf{PRISM-Style Prompting (B3)}~\citep{kirk2024prism},
(5)~\textbf{Activation Steering}~\citep{arditi2025refusal} using culturally
contrastive steering vectors (details in Appendix~\ref{app:baselines}).
Reward-guided methods (ARGS, controlled decoding) are excluded as they require
reward models unavailable for cross-cultural settings
(Appendix~\ref{app:baselines}).

\subsection{Metrics}
\label{sec:exp:metrics}

\textbf{JSD} (primary, $\downarrow$): Jensen-Shannon Distance between model and human
AMCE vectors. \textbf{Pearson $r$} ($\uparrow$): correlation of 6D AMCE vectors.
\textbf{MAE} ($\downarrow$): mean absolute AMCE error. All aggregated as unweighted
means across 15 countries with bootstrap 95\% CIs (1000 resamples).

% ============================================================
\section{Results}
\label{sec:results}
% ============================================================

\subsection{Cross-Model Comparison}

Table~\ref{tab:model_summary} summarises performance across six models. The two
70B-class models benefit most: Qwen2.5-72B achieves \textbf{37.6\%} mean JSD
improvement (Pearson $r$: $0.548 \to 0.696$) and Llama-3.1-70B achieves
\textbf{34.7\%} ($r$: $-0.026 \to 0.384$). Smaller models show consistent
+18--23\% gains. Mistral-Large (+5.8\%) has high baseline alignment leaving less
room for correction. Qwen2.5-32B ($-0.2\%$) exhibits concentrated logit distributions (low
inter-token entropy at the decision position) causing the fixed $\sigma=0.3$
perturbations to land outside the model's effective logit range, yielding
near-uniform importance weights ($K_{\text{eff}} \approx K$) and reducing
MPPI to random noise. A natural remedy is adaptive $\sigma$ selection: setting
$\sigma \propto \text{std}(\delta_1,\ldots,\delta_N)$ would scale perturbations
to each model's logit spread, ensuring meaningful importance-weight
differentiation. We leave empirical validation of this heuristic for future
work (Appendix~\ref{app:limitations}).

\begin{table}[t]
\caption{Cross-model comparison. Bootstrap 95\% CIs; differences $>$$\pm$0.004
significant at $p{<}0.05$.}
\label{tab:model_summary}
\centering\small
\setlength{\tabcolsep}{5pt}
\begin{tabular}{lcccccc}
\toprule
\multirow{2}{*}{Model} & \multirow{2}{*}{Size} &
  \multicolumn{2}{c}{Mean JSD $\downarrow$} &
  \multirow{2}{*}{\makecell{Improv.\\ (\%)}} &
  \multicolumn{2}{c}{Pearson $r$ $\uparrow$} \\
\cmidrule(lr){3-4}\cmidrule(lr){6-7}
 & & Van. & SWA & & Van. & SWA \\
\midrule
Qwen2.5-72B   & 72B  & $.0859 \pm .005$ & $.0536 \pm .003$ & \gain{\textbf{+37.6}} & 0.548 & \textbf{0.696} \\
Llama-3.1-70B & 70B  & $.0937 \pm .006$ & $.0612 \pm .004$ & \gain{+34.7}          & $-$.026 & 0.384 \\
Llama-3.1-8B  & 8B   & $.0533 \pm .004$ & $.0436 \pm .003$ & \gain{+22.6}          & $-$.247 & 0.271 \\
Qwen2.5-7B    & 7B   & $.0522 \pm .004$ & $.0426 \pm .003$ & \gain{+18.4}          & 0.252  & 0.402 \\
Mistral-Large & 126B & $.0457 \pm .003$ & $.0430 \pm .003$ & \gain{+5.8}           & 0.658  & 0.550 \\
Qwen2.5-32B   & 32B  & $.0632 \pm .004$ & $.0633 \pm .004$ & \loss{$-$0.2}         & 0.625  & 0.550 \\
\bottomrule
\end{tabular}
\end{table}

\subsection{Baseline Comparison on Llama-3.1-70B}

Table~\ref{tab:baselines}: SWA-MPPI achieves the lowest mean JSD (0.0612) and highest
Pearson $r$ (0.384), outperforming the best alternative (country-tailored prompting)
by 10.0\%. Activation steering performs \emph{below} vanilla, suggesting contrastive
steering vectors interfere with moral reasoning on this benchmark.

\begin{table}[t]
\caption{Baseline comparison on Llama-3.1-70B across 15 countries.}
\label{tab:baselines}
\centering\small
\begin{tabular}{lcccc}
\toprule
Method & Mean JSD $\downarrow$ & Pearson $r$ $\uparrow$ & $\Delta$JSD vs SWA \\
\midrule
\textbf{SWA-MPPI (ours)} & $\mathbf{.0612 \pm .004}$ & \textbf{0.384} & -- \\
Country-Tailored (B1) & $.0680 \pm .004$ & 0.233 & +11.1\% \\
PRISM-Style (B3)      & $.0699 \pm .005$ & 0.148 & +12.5\% \\
WVS Profile (B2)      & $.0737 \pm .005$ & 0.155 & +16.9\% \\
Activation Steering   & $.0877 \pm .006$ & $-$0.208 & +30.2\% \\
Vanilla               & $.0937 \pm .006$ & $-$0.026 & +34.7\% \\
\bottomrule
\end{tabular}
\end{table}

\subsection{Per-Country Results: Qwen2.5-72B}

Table~\ref{tab:results_qwen}: SWA-MPPI improves JSD in 13/15 countries.
\textbf{The Anglosphere paradox:} the largest gains occur for Australia (+75.4\%),
Great Britain (+72.2\%), and USA (+71.1\%)---precisely the English-speaking
populations one might expect to be best served. Their high vanilla JSD ($\approx$0.10--0.11) reflects RLHF-induced biases (extreme animal protection, strong
gender-neutral stances) that diverge from average human preferences.

\textbf{Germany (ceiling case):} vanilla JSD = 0.0603 with $r = 0.845$---already
exceptionally aligned. SWA-MPPI holds steady (JSD unchanged), demonstrating its
\emph{self-limiting} property: low inter-persona disagreement means the trigger
rarely fires. The Pearson $r$ decrease (0.845$\to$0.747) reflects sensitivity to
small rank perturbations in 6D, not distributional degradation
(Appendix~\ref{app:germany}).

Saudi Arabia ($-$2.3\%) and Brazil (+6.7\%) have limited WVS Wave~7 coverage,
constraining persona quality.

\begin{table}[t]
\caption{Per-country results for \textbf{Qwen2.5-72B-Instruct}.}
\label{tab:results_qwen}
\centering
\setlength{\tabcolsep}{4.5pt}
\footnotesize
\begin{tabular}{llcccccc}
\toprule
\multirow{2}{*}{Country} & \multirow{2}{*}{Lang.} &
  \multicolumn{2}{c}{JSD $\downarrow$} &
  \multicolumn{2}{c}{Pearson $r$ $\uparrow$} &
  \multicolumn{2}{c}{JSD Improv.} \\
\cmidrule(lr){3-4}\cmidrule(lr){5-6}
 & & Van. & SWA & Van. & SWA & ($\Delta$) & (\%) \\
\midrule
USA         & en & .1066 & \textbf{.0308} & 0.400 & 0.822 & $-$.076 & \gain{+71.1} \\
Germany     & de & .0603 & .0603          & 0.845 & 0.747 & +.000   & ~~+0.0 \\
China       & zh & .0513 & .0371          & 0.851 & 0.762 & $-$.014 & \gain{+27.8} \\
Japan       & ja & .0986 & .0553          & 0.634 & 0.786 & $-$.043 & \gain{+43.9} \\
Brazil      & pt & .0754 & .0703          & 0.349 & 0.529 & $-$.005 & \gain{+6.7}  \\
Saudi Arabia& ar & .0778 & .0796          & 0.443 & 0.527 & +.002   & \loss{$-$2.3}\\
Vietnam     & vi & .1302 & .0783          & 0.569 & 0.456 & $-$.052 & \gain{+39.9} \\
France      & fr & .0789 & .0479          & 0.617 & 0.838 & $-$.031 & \gain{+39.4} \\
India       & hi & .0600 & .0561          & 0.670 & 0.739 & $-$.004 & \gain{+6.5}  \\
South Korea & ko & .0720 & .0513          & 0.355 & 0.645 & $-$.021 & \gain{+28.8} \\
Great Britain& en& .1057 & \textbf{.0294} & 0.419 & 0.800 & $-$.076 & \gain{+72.2} \\
Russia      & ru & .0800 & .0629          & 0.623 & 0.784 & $-$.017 & \gain{+21.4} \\
Mexico      & es & .0760 & .0534          & 0.605 & 0.545 & $-$.023 & \gain{+29.8} \\
Nigeria     & en & .1103 & .0649          & 0.407 & 0.615 & $-$.046 & \gain{+41.2} \\
Australia   & en & .1047 & \textbf{.0258} & 0.438 & 0.841 & $-$.079 & \gain{+75.4} \\
\midrule
\textbf{Mean} & & \textbf{.0859} & \textbf{.0536} & 0.548 & 0.696 & $-$.032 & \gain{\textbf{+37.6}} \\
\bottomrule
\end{tabular}
\end{table}

% ============================================================
\section{Ablation Study}
\label{sec:ablation}
% ============================================================

We ablate six components on Qwen2.5-72B, USA (342 scenarios;
Table~\ref{tab:ablation}). Note: the ablation uses the full scenario pool, yielding
JSD = 0.0425 vs.\ 0.0308 in Table~\ref{tab:results_qwen} due to different dataset
splits (Appendix~\ref{app:ablation_detail}).

\begin{table}[t]
\caption{Ablation on Qwen2.5-72B, USA (342 scenarios). Full system:
JSD = 0.0425, $r$ = 0.724.}
\label{tab:ablation}
\centering\small
\setlength{\tabcolsep}{4pt}
\begin{tabular}{llccccc}
\toprule
\# & Configuration & JSD $\downarrow$ & $\Delta$JSD & $r$ $\uparrow$ & $\Delta r$ & MAE \\
\midrule
  & \textbf{Full SWA-MPPI} & \textbf{0.0425} & -- & \textbf{0.724} & -- & \textbf{9.77} \\
\midrule
1 & No-MPPI (consensus only)          & 0.0493 & \loss{+.007} & 0.650 & \loss{$-$.074} & 11.02 \\
2 & Always-on MPPI                    & 0.0422 & \gain{$-$.000} & 0.729 & \gain{+.005} & 9.74 \\
3a & Utility: quadratic               & 0.0429 & \loss{+.000} & 0.725 & +.001         & 9.68 \\
3b & Utility: linear                  & 0.0444 & \loss{+.002} & 0.706 & \loss{$-$.018} & 10.12 \\
4  & Without utilitarian ($N{=}3$)    & 0.0610 & \loss{+.019} & 0.441 & \loss{$-$.283} & 11.09 \\
5  & No debiasing                     & 0.0476 & \loss{+.005} & 0.743 & +.019         & 9.67 \\
\bottomrule
\end{tabular}
\end{table}

The results establish a clear hierarchy: (i)~the \textbf{utilitarian persona} is most
critical (+43.5\% JSD without it; Pearson $r$ collapses by 0.283), anchoring the
ensemble along the numerosity dimension; (ii)~\textbf{MPPI} adds meaningful
refinement (+15.9\% JSD); (iii)~\textbf{positional debiasing} corrects substantial
recency bias (+12.0\% JSD; Species preference inflated by 7.7pp without it);
(iv)~\textbf{Prospect Theory} contributes through nonlinear loss-averse weighting
(+4.5\% JSD vs.\ linear). The adaptive trigger fires at 100\% for USA---a genuine
empirical property of high-disagreement scenarios, not a calibration failure. The
MPPI \emph{flip rate} is only 6.0\%, indicating the framework \emph{modulates
confidence} rather than reversing decisions (Appendix~\ref{app:trigger}).

% ============================================================
\section{Discussion}
\label{sec:discussion}
% ============================================================

\paragraph{Key findings.}
SWA-MPPI achieves three desirable properties simultaneously: (i)~substantial
alignment gains on well-calibrated 70B models without weight modification,
(ii)~graceful self-limiting when intervention is unnecessary (Germany), and
(iii)~a 4.6$\times$ latency overhead (730ms vs.\ 160ms per scenario on H100)
negligible for batch inference. The MPPI loop operates on cached scalars and
contributes negligibly to wall-clock time.

\paragraph{The Anglosphere paradox.}
The largest gains occur for English-speaking countries---precisely those
best-represented in training data. This reveals that RLHF introduces specific moral
biases (extreme animal protection, strong gender-neutral stances) that \emph{actively
worsen} alignment with even Western human preferences, and SWA-MPPI's WVS-derived
personas successfully counterbalance these.

\paragraph{Cross-model failure patterns.}
Per-country results (Appendix~\ref{app:percountry_all}) reveal that failure modes
are \emph{model-specific}, not country-specific. On Llama-3.1-70B, the largest gains
shift to Vietnam (+65.8\%) and Mexico (+66.4\%), while Brazil ($-$23.3\%) and
Saudi Arabia ($-$10.7\%) regress---a different pattern from Qwen2.5-72B, where
Brazil improves (+6.7\%) and Saudi Arabia is the only regression ($-$2.3\%).
Mistral-Large shows a striking dissociation: JSD improves modestly (+5.8\%) but
Pearson~$r$ \emph{decreases} for 11/15 countries (mean $\Delta r = -0.108$),
indicating that the MPPI correction improves distributional distance while
perturbing the rank ordering of AMCE dimensions. Qwen2.5-32B degrades in
7/15 countries, confirming the logit-concentration failure mode described in
Section~\ref{sec:results}. These patterns suggest that SWA-MPPI's effectiveness
depends on the base model's logit calibration quality, and that a pre-flight
diagnostic (e.g., checking logit entropy) could predict when to activate the
framework.

\paragraph{Limitations.}
We evaluate 15 of 100+ MultiTP countries; extending to low-resource languages
remains untested. The per-agent reward $r_i$ centres corrections relative to the base
model, encoding an assumption that persona shifts approximate human
directions---validated at aggregate level but unverified per-dimension. The WVS
grounding encodes surveyed averages at a specific historical moment
(Wave~7, $\sim$2017--2022); temporal staleness and limited WVS coverage for some
countries constrain applicability. Extended limitations are in
Appendix~\ref{app:limitations}.

\paragraph{Broader impact and safeguards.}
SWA-MPPI could improve fairness and inclusivity by calibrating LLM moral behaviour
to diverse cultural preferences without costly retraining---the ``Anglosphere
paradox'' insight alone has implications for how alignment protocols are designed
and evaluated. However, the framework also risks aligning to harmful majority
preferences (e.g., discriminatory norms against minorities within a population)
or reifying stereotypes embedded in WVS survey averages. We identify three
safeguards for responsible deployment: (i)~\emph{governance oversight}:
SWA-MPPI should not override safety-critical model guardrails, and deployment
decisions should involve domain experts familiar with the target culture;
(ii)~\emph{periodic persona updates}: WVS data reflects a specific historical
moment (Wave~7, $\sim$2017--2022), and personas should be refreshed as newer
waves become available to prevent temporal staleness; (iii)~\emph{robustness
checks}: practitioners should audit per-dimension AMCE shifts to detect cases
where persona-driven corrections entrench harmful norms (e.g., extreme gender
biases in certain WVS profiles). The WVS grounding provides a degree of empirical
accountability---persona descriptions reflect surveyed population averages rather
than author preferences---but this does not eliminate the risk of culturally
legitimised harm.

% ============================================================
\section{Conclusion}
\label{sec:conclusion}
% ============================================================

We presented SWA-MPPI, a training-free inference-time framework for cross-cultural
moral alignment. By grounding persona construction in WVS data, exploiting
inter-persona disagreement as a trigger signal, and applying Prospect-Theory-weighted
MPPI in logit space, SWA-MPPI reduces JSD to human preferences by 34--38\% on 70B
models, outperforms all baselines, and gracefully refrains from correction when
alignment is already high. The Anglosphere paradox illustrates how RLHF biases worsen
alignment even for well-represented populations. Future directions include:
adaptive $\sigma$ selection based on per-model logit entropy for broader
compatibility; extension to multi-option preference settings; integration with
newer WVS waves to assess longitudinal drift in moral preferences;
disentangling language from culture effects via multilingual evaluation within
single countries; and developing safeguards against entrenching discriminatory
population-level norms---for instance, by flagging dimensions where
persona-driven corrections exceed predefined ethical bounds.

% ============================================================
\begin{ack}
We thank the creators of MultiTP~\citep{jin2025multitp} and the Moral Machine
experiment~\citep{awad2018moral} for public data. Experiments on NVIDIA H100 hardware.
\end{ack}

% ============================================================
\bibliographystyle{abbrvnat}
\bibliography{references}

% ============================================================
% APPENDIX
% ============================================================
\newpage
\appendix

\section{AMCE Estimation Details}
\label{app:amce}

\paragraph{Binary dimensions.} For Species, Gender, Age, Fitness, and Social Value,
each scenario presents a forced choice between two character groups differing on
exactly one dimension. The AMCE reduces to the empirical mean of
$p_{\text{spare}}(\mathbf{x})$ across all scenarios in that dimension
(Eq.~\ref{eq:mpr}), consistent with MultiTP conventions~\citep{jin2025multitp}.

\paragraph{Utilitarianism dimension.} Utilitarianism scenarios vary the \emph{number}
of lives on each side rather than character identity. Since the count difference
$n_{\text{diff}} = |n_{\text{pref}} - n_{\text{non-pref}}|$ varies across scenarios
(typically 1--3), a simple mean would conflate signal strength with preference
direction. We therefore fit an OLS regression
$p_{\text{spare}} = a + b \cdot n_{\text{diff}}$ and evaluate at the mean observed
$n_{\text{diff}}$ in the data. This avoids underestimation: evaluating at
$n_{\text{diff}}=1$ would undercount the utilitarian signal when typical scenarios
have $n_{\text{diff}} \geq 2$. Scenarios with $n_{\text{diff}}=0$ (equal group sizes)
are excluded as non-discriminative.

\paragraph{Human AMCE scale.} The MultiTP ground-truth AMCEs in
\texttt{country\_specific\_ACME.csv} are on a $[-1, 1]$ scale. We convert to
$[0, 100]$ via $(1 + \text{AMCE}_{\text{raw}}) / 2 \times 100$, where 50\%
represents no preference and 100\% represents maximal preference for the
``preferred'' group. Model AMCEs are computed on the same scale.

\paragraph{Consistency.} \citet{awad2018moral} show that marginal averages are
unbiased estimators under the balanced randomisation design of the Moral Machine,
making our estimation valid and reproducible.

\section{Persona Construction Details}
\label{app:personas}

\subsection{Example Personas}

The young-adult persona for the \textbf{United States}:
\begin{quote}
\small\itshape
You are a young adult (18--35) from the United States of America. Based on survey
data, people in your age group in the USA typically value: strong gender equality
(score: 8.2/10), low religious centrality in moral decisions (score: 3.1/10), high
personal autonomy (score: 7.8/10), high meritocratic orientation (score: 7.5/10),
moderate family importance (score: 6.4/10), and moderate interpersonal trust
(score: 5.3/10). When making moral judgments, you tend to emphasise individual
rights and equal treatment, and weigh personal achievement and rule-following.
\end{quote}

The older-adult persona for \textbf{Vietnam}:
\begin{quote}
\small\itshape
You are an older adult (55+) from Vietnam. Based on survey data, people in your age
group in Vietnam typically value: moderate gender equality (score: 5.1/10), high
religious and moral traditionalism (score: 7.2/10), strong family importance
(score: 9.4/10), moderate personal autonomy (score: 4.8/10), and low interpersonal
trust toward strangers (score: 3.2/10). When making moral judgments, you tend to
emphasise collective welfare, family obligations, and hierarchical social roles.
\end{quote}

The utilitarian persona (country-invariant): \textit{``You are a strict utilitarian.
Always choose the option that saves the greatest number of lives, regardless of the
characteristics of the individuals involved.''}

\subsection{WVS Data Processing Pipeline}
\label{app:wvs_pipeline}

We process WVS Wave~7 individual-level microdata as follows:
\begin{enumerate}
  \item \textbf{Variable extraction:} For each respondent, we extract 8 value
    dimension scores from specific WVS items: gender egalitarianism (mean of
    Q58P, Q59P, Q60P), religious importance (Q6P), interpersonal trust (Q43P),
    moral permissiveness (mean of Q50, Q52P, Q54P), work centrality (Q5P),
    family importance (Q1P), personal autonomy (Q39P), meritocratic orientation
    (Q40P). All items use inverted scales where higher = more
    progressive/positive.
  \item \textbf{Age cohort assignment:} Using birth year (Q261) and survey year
    (A\_YEAR), we compute age and assign to: \emph{young} ($<$36),
    \emph{middle} (36--55), \emph{older} ($>$55). Respondents with birth year
    before 1900 or after 2010, or survey year before 2015, are excluded.
  \item \textbf{Cohort aggregation:} For each country and age cohort, we compute
    the mean of each dimension across all valid respondents.
  \item \textbf{Natural-language conversion:} Dimension means are converted to
    descriptors using calibrated thresholds (e.g., religiosity ratio $>$0.85
    $\to$ ``deeply religious''; trust ratio $>$0.55 $\to$ ``high interpersonal
    trust'').
  \item \textbf{Fallback:} For countries with insufficient WVS coverage (e.g.,
    Saudi Arabia), we use manually authored native-language personas grounded in
    area-studies literature, covering the same demographic structure (3 age
    cohorts + 1 utilitarian).
\end{enumerate}

\subsection{WVS-to-Trolley Dimension Linkage}
\label{app:wvs_linkage}

The linkage between WVS value dimensions and trolley moral dimensions operates
through two mechanisms.

\emph{Direct thematic overlap:} WVS \emph{gender egalitarianism} maps onto the
Gender dimension (societies scoring higher show weaker female-sparing
preferences~\citep{awad2018moral}); \emph{religious importance} correlates with
Species (human exceptionalism); \emph{meritocratic orientation} relates to Social
Value and Fitness.

\emph{Indirect modulation:} Dimensions like \emph{personal autonomy},
\emph{interpersonal trust}, and \emph{family importance} modulate the moral
\emph{frame}---e.g., high-autonomy personas favour utilitarian calculi, while
high-family-importance personas weight relational roles (affecting Age and Social
Value). Crucially, SWA-MPPI does not require exact causal mapping: the ensemble
generates diversity in logit-space gaps, and MPPI selects the correction balancing
collective utility. A rigorous per-dimension causal analysis (e.g., regressing WVS
attributes against AMCE shifts) would further strengthen this argument.

\section{MPPI Theoretical Justification}
\label{app:mppi_theory}

\subsection{Degenerate but Principled Stochastic Optimal Control}

The scalar decision gap $\delta$ constitutes a one-dimensional \emph{state};
``dynamics'' are identity (no state transition); perturbations $\epsilon_k$ are
``control inputs'' at a single horizon step. This is a degenerate case of
MPPI~\citep{williams2017mppi}, but the information-theoretic derivation remains
valid: $w_k \propto \exp(U_{\text{total}}/\beta)$ minimises KL divergence between
the sample distribution and the Boltzmann distribution of utility, making the KL
term in Eq.~\ref{eq:util_total} a direct consequence of the free-energy
objective---not a heuristic penalty.

\subsection{Comparison with Simpler Alternatives}

Temperature scaling uniformly stretches the logit gap without accounting for
direction or magnitude of persona rewards. A fixed margin offset ignores the
utility landscape, treating all perturbations as equally valuable. MPPI performs an
\emph{adaptive, scenario-specific} correction via nonlinear importance weighting.
The ablation (Table~\ref{tab:ablation}) provides indirect evidence: replacing PT
with linear utility (closer to uniform weighting, hence a simple mean shift)
degrades JSD by 4.5\%.

\section{Hyperparameter Validation}
\label{app:hyperparams}

Non-PT hyperparameters ($\sigma$, $\lambda_{\text{coop}}$, $\alpha_{\text{KL}}$,
$T_{\text{dec}}$, $T_{\text{cat}}$) were validated on a held-out synthetic set of
200 binary moral dilemmas generated by GPT-4, covering the same six dimensions as
MultiTP with novel character descriptions. Pseudo-ground-truth labels from three
annotators (the authors) provided proxy AMCEs. Grid search minimised mean JSD on
Qwen2.5-72B. Transferability was verified on a disjoint 50-scenario MultiTP English
subset: performance within 1.2\% relative JSD of oracle-tuned configuration.
$T_{\text{cat}}$ values reflect empirical logit magnitude distributions: Species
and Gender produce systematically larger gaps, requiring higher temperatures.

\section{Dataset Preprocessing and Calibration}
\label{app:dataset}
\label{app:calibration}

\subsection{Preprocessing Sensitivity}

Up-sampling duplicates scenarios without modification, preserving AMCE signal while
reducing variance. Category capping prevents dimension dominance. Side balancing
enables effective debiasing. A formal sensitivity analysis (varying capping ratios,
up-sampling rates) is recommended for future work.

\subsection{Calibration Limitation}

The 50-scenario $\tau^{(c)}$ calibration set is drawn from the evaluation
distribution. We mitigate overfitting risk: (i)~calibration observes only a single
scalar (variance), unlikely to overfit; (ii)~it targets a fixed trigger rate, not
alignment metrics; (iii)~the 100\% trigger rate for USA suggests $\tau$ is not
sensitively tuned. A fully separate calibration corpus would provide stronger
guarantees.

\section{Cross-Lingual Token Elicitation}
\label{app:token_elicitation}

\paragraph{Native-language prompting system.}
Each scenario is rendered entirely in the target country's native language.
The implementation includes full translations of: (i)~scenario framing and
context sentences (4 paraphrase variants per language), (ii)~all 22 character
types with singular/plural forms in each language, (iii)~lane labels and group
identifiers in native script, and (iv)~closing questions with
language-appropriate conjunctions and grammar.

The prompt frame wraps the scenario in a native-language instruction that
explicitly requests an English answer token. Each frame follows the pattern:
``[moral dilemma preamble in native language] \textbackslash n \{scenario\}
\textbackslash n [instruction to answer LEFT or RIGHT in English]''.
This ensures moral reasoning occurs in the culturally native linguistic
frame while maintaining a consistent decision interface.

\paragraph{Token verification.}
Token IDs are verified per model at initialisation by encoding the uppercase
strings \texttt{"LEFT"} and \texttt{"RIGHT"} with
\texttt{add\_special\_tokens=False} and asserting single-token mappings:
Qwen2.5: LEFT=22397, RIGHT=11572; Llama-3.1: LEFT=31266, RIGHT=9268.
Logits for these two positions are extracted from the last-position output
tensor; no generation sampling is performed. The model uses its built-in
\texttt{chat\_template} for correct role formatting across model families.

\section{Baseline Implementation Details}
\label{app:baselines}

\subsection{Activation Steering}
For each country, a steering vector is computed as the difference in mean
residual-stream activations from 20 culturally contrastive prompt pairs (e.g.,
``Answer as someone from [country] who values [WVS-high trait]'' vs.\ ``Answer
as someone with the opposite values''). Vectors are extracted from layer~32
(transformer midpoint) of Llama-3.1-70B and applied at inference time by adding
$\alpha \cdot \mathbf{v}_{\text{steer}}$ to the residual stream, with
$\alpha=1.5$ selected via the same synthetic validation set used for SWA-MPPI.

\paragraph{Sensitivity analysis and fairness considerations.}
Activation steering performance is known to be sensitive to (i)~the extraction
layer, (ii)~the scaling coefficient $\alpha$, and (iii)~the construction of
contrastive pairs~\citep{arditi2025refusal}. We conducted a limited sensitivity
check: extracting from layers 24, 32, and 40 with $\alpha \in \{0.5, 1.0, 1.5,
2.0, 3.0\}$, we found that layer~32 with $\alpha=1.5$ yielded the lowest mean
JSD (0.0877); other configurations ranged from 0.089 to 0.112. We acknowledge
that a more exhaustive search---varying contrastive prompt design, using
multiple layers simultaneously, or employing PCA-based direction
refinement---could improve activation steering performance. However, even the
best configuration performs below vanilla on this benchmark, suggesting a
fundamental mismatch between linear activation interventions and the
multi-dimensional moral calibration problem. The key structural advantage of
SWA-MPPI is that it operates \emph{per-scenario} with adaptive correction
magnitude, whereas activation steering applies a fixed direction uniformly
across all scenarios.

\subsection{Omitted Baselines}
ARGS~\citep{khanov2024args} and controlled decoding~\citep{mudgal2024controlled}
require trained reward models. In the cross-cultural setting, this demands either a
separate reward model per country or culturally conditioned models---neither exists
for MultiTP. SWA-MPPI's constraint is operating without labelled preference data,
making these methodologically incompatible.

\section{Detailed Ablation Analysis}
\label{app:ablation_detail}

\textbf{JSD difference with Table~\ref{tab:results_qwen}:} The ablation uses the
full 342-scenario pool; per-country results use the evaluation split after
calibration reservation, yielding different scenario compositions.

\textbf{Always-on vs.\ adaptive trigger:} For USA, the adaptive trigger fires at
100\%---both strategies are numerically equivalent. The adaptive mechanism provides
savings in lower-disagreement countries.

\textbf{Utilitarian persona:} Removing it collapses Utilitarianism MPR from 61.2\%
to 48.0\% (human: 76.6\%), confirming it anchors the numerosity dimension.

\textbf{Debiasing:} Mean positional bias is 0.193 logit units; Species preference
inflates from 85.8\% to 93.5\% without it (human: 79.2\%).

\section{Trigger Mechanism and Latency}
\label{app:trigger}

The adaptive $\tau$ targets 35\% activation but achieves 100\% for USA (all scenarios
contested). The \emph{flip rate} (MPPI reversals) is 6.0\%, indicating confidence
modulation rather than decision reversal.

Mean latency: 730ms (SWA-MPPI) vs.\ 160ms (vanilla) on Qwen2.5-72B (H100).
Overhead from batched multi-agent passes and debiasing; MPPI operates on cached
scalars and contributes negligibly.

\section{Germany Pearson $r$ Analysis}
\label{app:germany}

Germany's JSD remains unchanged (0.0603) while Pearson $r$ drops from 0.845 to 0.747.
In 6D, Pearson $r$ is sensitive to rank perturbations: a single dimension swap reduces
$r$ by $\sim$0.1 while JSD (depending on magnitudes, not ordering) stays constant.
Rare MPPI corrections introduce minor rank perturbations without affecting
distributional distance.

\section{Per-Country Results for All Models}
\label{app:percountry_all}

Tables~\ref{tab:llama70b_percountry}--\ref{tab:qwen32b_percountry} present
per-country results for all evaluated models beyond Qwen2.5-72B
(Table~\ref{tab:results_qwen} in the main text).

\begin{table}[h]
\caption{Per-country results for \textbf{Llama-3.1-70B-Instruct}.}
\label{tab:llama70b_percountry}
\centering\setlength{\tabcolsep}{3.5pt}\footnotesize
\begin{tabular}{llcccccc}
\toprule
Country & Lang. & \multicolumn{2}{c}{JSD $\downarrow$} & \multicolumn{2}{c}{Pearson $r$ $\uparrow$} & \multicolumn{2}{c}{JSD Improv.} \\
\cmidrule(lr){3-4}\cmidrule(lr){5-6}
 & & Van. & SWA & Van. & SWA & ($\Delta$) & (\%) \\
\midrule
USA         & en & .1561 & .1234 & 0.012 & 0.197 & $-$.033 & \gain{+20.9} \\
Germany     & de & .0709 & .0391 & $-$0.455 & 0.685 & $-$.032 & \gain{+44.9} \\
China       & zh & .0459 & .0311 & 0.406 & 0.802 & $-$.015 & \gain{+32.3} \\
Japan       & ja & .0457 & .0416 & 0.496 & 0.466 & $-$.004 & \gain{+9.0} \\
Brazil      & pt & .0532 & .0656 & $-$0.370 & 0.348 & +.012  & \loss{$-$23.3} \\
Saudi Arabia& ar & .0539 & .0596 & 0.043 & $-$0.210 & +.006 & \loss{$-$10.7} \\
Vietnam     & vi & .1546 & .0529 & $-$0.519 & $-$0.135 & $-$.102 & \gain{+65.8} \\
France      & fr & .0500 & .0516 & 0.250 & 0.394 & +.002  & \loss{$-$3.0} \\
India       & hi & .0442 & .0337 & $-$0.315 & 0.715 & $-$.010 & \gain{+23.6} \\
South Korea & ko & .0893 & .0436 & $-$0.272 & 0.551 & $-$.046 & \gain{+51.2} \\
Great Britain& en& .1568 & .0918 & 0.012 & 0.286 & $-$.065 & \gain{+41.5} \\
Russia      & ru & .0449 & .0413 & 0.539 & 0.637 & $-$.004 & \gain{+7.9} \\
Mexico      & es & .1265 & .0425 & $-$0.371 & 0.307 & $-$.084 & \gain{+66.4} \\
Nigeria     & en & .1582 & .0988 & 0.109 & 0.497 & $-$.059 & \gain{+37.5} \\
Australia   & en & .1551 & .1010 & 0.039 & 0.224 & $-$.054 & \gain{+34.9} \\
\midrule
\textbf{Mean} & & \textbf{.0937} & \textbf{.0612} & $-$0.026 & 0.384 & $-$.033 & \gain{\textbf{+34.7}} \\
\bottomrule
\end{tabular}
\end{table}

\begin{table}[h]
\caption{Per-country results for \textbf{Qwen2.5-7B-Instruct} (4-bit).}
\label{tab:qwen7b_percountry}
\centering\setlength{\tabcolsep}{3.5pt}\footnotesize
\begin{tabular}{llcccccc}
\toprule
Country & Lang. & \multicolumn{2}{c}{JSD $\downarrow$} & \multicolumn{2}{c}{Pearson $r$} & \multicolumn{2}{c}{JSD Improv.} \\
\cmidrule(lr){3-4}\cmidrule(lr){5-6}
 & & Van. & SWA & Van. & SWA & ($\Delta$) & (\%) \\
\midrule
USA         & en & .0531 & .0448 & 0.048 & 0.531 & $-$.008 & \gain{+15.5} \\
Germany     & de & .0530 & .0522 & 0.207 & 0.212 & $-$.001 & \gain{+1.4} \\
China       & zh & .0512 & .0455 & 0.401 & 0.397 & $-$.006 & \gain{+11.2} \\
Japan       & ja & .0527 & .0540 & 0.490 & 0.449 & +.001 & \loss{$-$2.3} \\
Brazil      & pt & .0445 & .0331 & 0.415 & 0.837 & $-$.011 & \gain{+25.8} \\
Saudi Arabia& ar & .0539 & .0443 & 0.234 & 0.265 & $-$.010 & \gain{+17.9} \\
Vietnam     & vi & .0547 & .0472 & 0.259 & 0.263 & $-$.008 & \gain{+13.6} \\
France      & fr & .0529 & .0442 & 0.259 & 0.315 & $-$.009 & \gain{+16.5} \\
India       & hi & .0570 & .0524 & 0.466 & 0.183 & $-$.005 & \gain{+8.1} \\
South Korea & ko & .0492 & .0329 & 0.239 & 0.363 & $-$.016 & \gain{+33.2} \\
Great Britain& en& .0520 & .0302 & 0.140 & 0.483 & $-$.022 & \gain{+41.9} \\
Russia      & ru & .0506 & .0452 & 0.373 & 0.491 & $-$.005 & \gain{+10.5} \\
Mexico      & es & .0479 & .0384 & 0.038 & 0.360 & $-$.010 & \gain{+19.9} \\
Nigeria     & en & .0551 & .0411 & 0.101 & 0.425 & $-$.014 & \gain{+25.3} \\
Australia   & en & .0550 & .0431 & 0.170 & 0.452 & $-$.012 & \gain{+21.7} \\
\midrule
\textbf{Mean} & & \textbf{.0522} & \textbf{.0426} & 0.252 & 0.402 & $-$.010 & \gain{\textbf{+18.4}} \\
\bottomrule
\end{tabular}
\end{table}

\begin{table}[h]
\caption{Per-country results for \textbf{Llama-3.1-8B-Instruct} (4-bit).}
\label{tab:llama8b_percountry}
\centering\setlength{\tabcolsep}{3.5pt}\footnotesize
\begin{tabular}{llcccccc}
\toprule
Country & Lang. & \multicolumn{2}{c}{JSD $\downarrow$} & \multicolumn{2}{c}{Pearson $r$} & \multicolumn{2}{c}{JSD Improv.} \\
\cmidrule(lr){3-4}\cmidrule(lr){5-6}
 & & Van. & SWA & Van. & SWA & ($\Delta$) & (\%) \\
\midrule
USA         & en & .0470 & .0333 & $-$0.323 & 0.498 & $-$.014 & \gain{+29.1} \\
Germany     & de & .0716 & .0529 & $-$0.101 & 0.382 & $-$.019 & \gain{+26.1} \\
China       & zh & .0561 & .0457 & $-$0.082 & 0.388 & $-$.010 & \gain{+18.5} \\
Japan       & ja & .0518 & .0519 & 0.545 & 0.395 & +.000 & \loss{$-$0.2} \\
Brazil      & pt & .0732 & .0628 & $-$0.222 & 0.109 & $-$.010 & \gain{+14.2} \\
Saudi Arabia& ar & .0551 & .0513 & $-$0.252 & $-$0.230 & $-$.004 & \gain{+6.9} \\
Vietnam     & vi & .0505 & .0423 & $-$0.330 & 0.303 & $-$.008 & \gain{+16.2} \\
France      & fr & .0515 & .0612 & $-$0.610 & $-$0.220 & +.010 & \loss{$-$18.8} \\
India       & hi & .0756 & .0564 & $-$0.471 & 0.347 & $-$.019 & \gain{+25.3} \\
South Korea & ko & .0522 & .0419 & $-$0.352 & $-$0.062 & $-$.010 & \gain{+19.7} \\
Great Britain& en& .0375 & .0255 & $-$0.065 & 0.457 & $-$.012 & \gain{+32.0} \\
Russia      & ru & .0560 & .0434 & $-$0.230 & 0.593 & $-$.013 & \gain{+22.5} \\
Mexico      & es & .0439 & .0383 & $-$0.079 & 0.344 & $-$.006 & \gain{+12.7} \\
Nigeria     & en & .0474 & .0410 & $-$0.229 & 0.303 & $-$.006 & \gain{+13.5} \\
Australia   & en & .0505 & .0363 & $-$0.156 & 0.328 & $-$.014 & \gain{+28.1} \\
\midrule
\textbf{Mean} & & \textbf{.0533} & \textbf{.0436} & $-$0.247 & 0.271 & $-$.011 & \gain{\textbf{+22.6}} \\
\bottomrule
\end{tabular}
\end{table}

\begin{table}[h]
\caption{Per-country results for \textbf{Mistral-Large-2407} (4-bit).
Note: despite modest JSD gains, Pearson~$r$ decreases for 11/15 countries.}
\label{tab:mistral_percountry}
\centering\setlength{\tabcolsep}{3.5pt}\footnotesize
\begin{tabular}{llcccccc}
\toprule
Country & Lang. & \multicolumn{2}{c}{JSD $\downarrow$} & \multicolumn{2}{c}{Pearson $r$} & \multicolumn{2}{c}{JSD Improv.} \\
\cmidrule(lr){3-4}\cmidrule(lr){5-6}
 & & Van. & SWA & Van. & SWA & ($\Delta$) & (\%) \\
\midrule
USA         & en & .0414 & .0211 & 0.841 & 0.834 & $-$.020 & \gain{+49.0} \\
Germany     & de & .0378 & .0291 & 0.734 & 0.675 & $-$.009 & \gain{+23.0} \\
China       & zh & .0572 & .0618 & 0.824 & 0.585 & +.005 & \loss{$-$8.1} \\
Japan       & ja & .0457 & .0535 & 0.814 & 0.387 & +.008 & \loss{$-$17.1} \\
Brazil      & pt & .0519 & .0625 & 0.359 & 0.510 & +.011 & \loss{$-$20.4} \\
Saudi Arabia& ar & .0238 & .0315 & 0.724 & 0.671 & +.008 & \loss{$-$32.4} \\
Vietnam     & vi & .0521 & .0428 & 0.768 & 0.571 & $-$.009 & \gain{+17.8} \\
France      & fr & .0481 & .0475 & 0.683 & 0.628 & $-$.001 & \gain{+1.2} \\
India       & hi & .0457 & .0527 & 0.772 & 0.631 & +.007 & \loss{$-$15.4} \\
South Korea & ko & .0332 & .0320 & 0.546 & 0.355 & $-$.001 & \gain{+3.6} \\
Great Britain& en& .0418 & .0331 & 0.853 & 0.730 & $-$.009 & \gain{+20.7} \\
Russia      & ru & .0552 & .0566 & 0.768 & 0.631 & +.001 & \loss{$-$2.6} \\
Mexico      & es & .0472 & .0481 & 0.748 & 0.642 & +.001 & \loss{$-$1.9} \\
Nigeria     & en & .0401 & .0301 & 0.128 & 0.131 & $-$.010 & \gain{+25.0} \\
Australia   & en & .0516 & .0338 & 0.865 & 0.682 & $-$.018 & \gain{+34.5} \\
\midrule
\textbf{Mean} & & \textbf{.0457} & \textbf{.0430} & 0.658 & 0.550 & $-$.003 & \gain{\textbf{+5.8}} \\
\bottomrule
\end{tabular}
\end{table}

\begin{table}[h]
\caption{Per-country results for \textbf{Qwen2.5-32B-Instruct} (4-bit).
SWA-MPPI degrades in 7/15 countries due to concentrated logit distributions.}
\label{tab:qwen32b_percountry}
\centering\setlength{\tabcolsep}{3.5pt}\footnotesize
\begin{tabular}{llcccccc}
\toprule
Country & Lang. & \multicolumn{2}{c}{JSD $\downarrow$} & \multicolumn{2}{c}{Pearson $r$} & \multicolumn{2}{c}{JSD Improv.} \\
\cmidrule(lr){3-4}\cmidrule(lr){5-6}
 & & Van. & SWA & Van. & SWA & ($\Delta$) & (\%) \\
\midrule
USA         & en & .0447 & .0440 & 0.887 & 0.738 & $-$.001 & \gain{+1.5} \\
Germany     & de & .0587 & .0534 & 0.893 & 0.803 & $-$.005 & \gain{+9.1} \\
China       & zh & .0507 & .0574 & 0.827 & 0.818 & +.007 & \loss{$-$13.1} \\
Japan       & ja & .0485 & .0592 & 0.846 & 0.805 & +.011 & \loss{$-$22.0} \\
Brazil      & pt & .0717 & .0645 & 0.552 & 0.422 & $-$.007 & \gain{+9.7} \\
Saudi Arabia& ar & .0883 & .1015 & 0.704 & 0.625 & +.013 & \loss{$-$14.9} \\
Vietnam     & vi & .0900 & .0780 & 0.634 & 0.556 & $-$.012 & \gain{+13.1} \\
France      & fr & .0730 & .0681 & 0.647 & 0.652 & $-$.005 & \gain{+6.7} \\
India       & hi & .0737 & .0840 & 0.458 & 0.583 & +.010 & \loss{$-$13.8} \\
South Korea & ko & .0529 & .0303 & 0.106 & 0.064 & $-$.023 & \gain{+43.0} \\
Great Britain& en& .0470 & .0345 & 0.889 & 0.732 & $-$.013 & \gain{+26.5} \\
Russia      & ru & .0700 & .0780 & 0.721 & 0.711 & +.008 & \loss{$-$11.3} \\
Mexico      & es & .0775 & .0933 & 0.632 & 0.381 & +.016 & \loss{$-$20.3} \\
Nigeria     & en & .0641 & .0512 & 0.124 & 0.156 & $-$.013 & \gain{+20.1} \\
Australia   & en & .0541 & .0425 & 0.883 & 0.781 & $-$.012 & \gain{+21.4} \\
\midrule
\textbf{Mean} & & \textbf{.0632} & \textbf{.0633} & 0.625 & 0.550 & +.000 & \loss{\textbf{$-$0.2}} \\
\bottomrule
\end{tabular}
\end{table}

\section{Extended Limitations}
\label{app:limitations}

\paragraph{Inter-persona reward assumption.}
The reward $r_i = \delta_i - \delta_{\text{base}}$ assumes persona shifts approximate
human target directions. This is plausible when the base model represents a WEIRD-biased
prior, but if a persona's shift is orthogonal to the human target, the signal may
mislead. Per-dimension validation is needed.

\paragraph{Country coverage.}
15 countries across major cultural clusters, but no low-resource languages (Amharic,
Khmer, Yoruba). Cross-lingual generalisation to such settings remains untested.

\paragraph{Binary format.}
The logit-difference formulation requires generalisation to softmax for multi-option
settings.

\paragraph{Language--culture conflation.}
One language per country may conflate linguistic and cultural effects.

\paragraph{Model quantisation.}
4-bit/Int8 quantisation may introduce logit artefacts affecting MPPI stability.

\paragraph{WVS temporal staleness.}
Wave~7 ($\sim$2017--2022) data may not reflect current preferences in rapidly
evolving societies.

\paragraph{Qwen2.5-32B failure mode.}
Concentrated logit distributions cause MPPI perturbations to be off-manifold,
yielding near-uniform importance weights. Adaptive $\sigma$ based on per-model
logit entropy could address this.

\end{document}
